\chapter{Resultados y evaluación}
\label{cap:resultadosEvaluacion}

En este capítulo se recogen los resultados obtenidos después de integrar las distintas partes del
proyecto. No se pretende repetir aquí todo el proceso de desarrollo, descrito en el capítulo
\ref{cap:descripcionTrabajo}, sino analizar qué se ha conseguido finalmente y hasta qué punto los
resultados cumplen los objetivos planteados al inicio del trabajo.

En primer lugar, se comparan los recursos hardware
utilizados por una misma operación implementada mediante \ac{HLS} y mediante \ac{VHDL}.
Después se analiza el
encaje del acelerador \ac{CIFAR}-10 dentro de la Zybo Z-7010. Por último, se estudia el tiempo de
inferencia obtenido en \ac{FPGA} frente a la ejecución del \textit{core} \texttt{C++}
generado por hls4ml en el \ac{ARM}.

\section{Comparación de recursos entre HLS y VHDL}

Antes de abordar el modelo \ac{CIFAR}-10 completo, se utilizó la multiplicación de
matrices de la fase 3
como prueba controlada. Esta prueba era más sencilla que la red neuronal, pero ya permitía estudiar
una cuestión importante: si Vitis \ac{HLS} generaba hardware razonable frente a una
descripción manual
en \ac{VHDL}.

Los recursos obtenidos para ambas implementaciones se muestran en la
tabla~\ref{tab:recursosMatricesHlsVhdl}.

\begin{table}[h!]
  \centering
  \begin{tabular}{|l|r|r|r|r|}
    \hline
    \textbf{Diseño} & \textbf{\ac{LUT}} & \textbf{FF} & \textbf{\ac{BRAM}} &
    \textbf{\ac{DSP}} \\ \hline
    \textbf{\ac{VHDL}} & 8263 & 4363 & 35 & 72 \\ \hline
    \textbf{\ac{HLS}}  & 7723 & 6894 & 32 & 72 \\ \hline
  \end{tabular}
  \caption{Comparación de recursos entre la multiplicación de matrices en VHDL y en HLS}
  \label{tab:recursosMatricesHlsVhdl}
\end{table}

Ambas versiones utilizan 72 \ac{DSP}, esto era esperable, ya que las dos
implementaciones explotan el mismo tipo de paralelismo en la parte aritmética, pero nos asegura
que ambas soluciones están alineadas. La diferencia está
en cómo se organiza el resto del circuito. La versión generada con \ac{HLS} usa 540
\ac{LUT} menos y 3 \ac{BRAM}
menos que la versión \ac{VHDL}, pero a cambio utiliza 2531 \ac{FF} más.

Este resultado fue útil por dos motivos. Primero, porque confirmó que \ac{HLS} no estaba
generando una
implementación desproporcionada en área. Segundo, porque mostró una de las ventajas prácticas de la
herramienta: el reparto entre memoria, registros y lógica de control se consigue con menos esfuerzo
manual.

Durante el desarrollo de la versión \ac{VHDL} apareció una primera implementación en la que no
se inferían correctamente las \ac{BRAM} y el uso de \ac{LUT} se disparaba. En \ac{HLS}
este tipo de decisiones
quedan más guiadas por la herramienta, lo cual no elimina la necesidad de revisar informes, pero sí
reduce mucho el coste de llegar a una primera versión funcional.

\section{Utilización de recursos del acelerador CIFAR-10}

El resultado más delicado del proyecto fue conseguir que el modelo \ac{CIFAR}-10 entrase en la Zybo
Legacy. La Zynq Z-7010 integra una \ac{FPGA} pequeña para implementar una red
convolucional, por lo que el
diseño estuvo limitado principalmente por el número de \ac{LUT} disponibles. A lo largo
del desarrollo
se probaron varias configuraciones de modelo y de cuantización hasta encontrar una versión que
Vivado pudiera colocar y rutear correctamente.

La versión final utilizada en placa fue el \ac{IP} generado a partir del modelo:

\texttt{v4\_8\_14\_32\_24\_final}

Tras integrarlo en Vivado, el diseño completo obtuvo la
utilización mostrada en la tabla~\ref{tab:recursosCifar10Final}.

\begin{table}[h!]
  \centering
  \begin{tabular}{|l|r|r|}
    \hline
    \textbf{Recurso} & \textbf{Uso en el diseño final} & \textbf{Porcentaje de uso} \\ \hline
    \ac{LUT} &   16117 & 91.57\% \\ \hline
    \ac{FF} &    17917 & 50.90\% \\ \hline
    \ac{BRAM} &  37,5  & 62.50\% \\ \hline
    \ac{DSP} &   6     & 7.50\%  \\ \hline
  \end{tabular}
  \caption{Utilización de recursos del diseño CIFAR-10 final en la Zybo}
  \label{tab:recursosCifar10Final}
\end{table}

El diseño queda cerca del límite de \ac{LUT}, pero entra en la placa. Esto es importante
porque varias
versiones anteriores fallaban en Vivado con errores de sobreutilización o de empaquetado
de \textit{slices}.
En cambio, los \ac{FF} no llegaron a ser el recurso crítico y el uso de \ac{BRAM} se
mantuvo dentro de lo
admisible para la placa.

Un aspecto llamativo es el bajo uso de \ac{DSP}
\citep{amdVitisHLS}. Durante el proyecto se intentó aumentar su utilización
para descargar lógica combinacional, pero no siempre es posible trasladar automáticamente el coste
de \ac{LUT} a \ac{DSP}. La decisión depende de los tipos de datos, de los anchos de palabra, de las
directivas \ac{HLS}, del factor de reutilización y de cómo Vivado termina mapeando las
operaciones. En
este caso, el límite real lo marcó la lógica auxiliar y de control asociada a la red, no solo las
multiplicaciones.

\section{Validación funcional de la inferencia}

Una vez generado el \textit{bitstream}, se realizaron pruebas con varias imágenes de prueba.

La primera ejecución se hizo sobre imágenes de test que entrega el propio \textit{dataset}
de \ac{CIFAR}-10, las cuales se han tenido que transformar al formato Q16.16 normalizado.

La ejecución sobre todas estas imágenes ha sido la siguiente:

\begin{table}[h!]
  \centering
  \resizebox{\textwidth}{!}{
    \begin{tabular}{|c|l|r|r|r|r|r|r|r|r|r|r|}
      \hline
      \textbf{Img. enviada} &
      \textbf{Pred.} &
      \textbf{Plane} &
      \textbf{Auto} &
      \textbf{Bird} &
      \textbf{Cat} &
      \textbf{Deer} &
      \textbf{Dog} &
      \textbf{Frog} &
      \textbf{Horse} &
      \textbf{Ship} &
      \textbf{Truck} \\ \hline

      Plane & \textcolor{green!60!black}{Plane} & \textbf{71.17\%} & 1.35\% & 11.78\% &
      0.60\% & 0.75\% & 0.22\% & 0.28\% & 0.03\% & 13.61\% & 0.21\% \\ \hline
      Auto & \textcolor{green!60!black}{Auto}   & 3.59\% & \textbf{60.70\%} & 0.26\% &
      2.88\% & 0.01\% & 10.63\% & 0.32\% & 6.72\% & 0.02\% & 14.88\% \\ \hline
      Bird & \textcolor{red}{Truck}             & 5.21\% & 2.07\% & 21.93\% & 6.79\% &
      19.49\% & 9.05\% & 2.06\% & 1.80\% & 1.35\% & \textbf{30.24\%} \\ \hline
      Cat & \textcolor{green!60!black}{Cat}     & 0.53\% & 0.32\% & 1.19\% &
      \textbf{53.14\%} & 2.07\% & 15.78\% & 26.00\% & 0.22\% & 0.45\% & 0.30\% \\ \hline
      Deer & \textcolor{green!60!black}{Deer}   & 32.07\% & 0.54\% & 15.93\% & 2.48\% &
      \textbf{34.68\%} & 2.32\% & 0.93\% & 1.53\% & 8.30\% & 1.22\% \\ \hline
      Dog & \textcolor{red}{Frog}               & 0.43\% & 1.65\% & 4.79\% & 9.69\% &
      17.14\% & 18.70\% & \textbf{32.13\%} & 14.99\% & 0.19\% & 0.28\% \\ \hline
      Frog & \textcolor{green!60!black}{Frog}   & 0.06\% & 0.13\% & 6.18\% & 5.56\% &
      3.72\% & 1.17\% & \textbf{83.11\%} & 0.02\% & 0.04\% & 0.01\% \\ \hline
      Horse & \textcolor{red}{Truck}            & 0.04\% & 20.13\% & 0.04\% & 0.30\% &
      0.48\% & 5.63\% & 0.80\% & 35.43\% & 0.02\% & \textbf{37.13\%} \\ \hline
      Ship & \textcolor{red}{Auto}              & 14.65\% & \textbf{71.26\%} & 0.25\% &
      0.02\% & 0.03\% & 0.01\% & 0.09\% & 0.01\% & 12.09\% & 1.60\% \\ \hline
      Truck & \textcolor{green!60!black}{Truck} & 0.17\% & 9.78\% & 0.06\% & 0.11\% &
      0.02\% & 0.35\% & 0.12\% & 0.07\% & 0.24\% & \textbf{89.08\%} \\ \hline

    \end{tabular}
  }
  \caption{Resultados de inferencia para 10 imágenes CIFAR-10 mostrando la confianza por clase}
  \label{tab:resultadosInferenciaCifar10}
\end{table}

Con estos resultados, podemos decir que la inferencia en la Zybo funciona de forma real,
aunque el modelo no es robusto.
Teniendo en cuenta que el modelo ha sido optimizado para que entre de manera ajustada en los
recursos que tiene la Z-7010, los resultados son aceptables para el estudio de este proyecto y
poder verificar que en efecto, podemos ejecutar inferencia.

Muchos de los fallos de predicción que vemos, están casi empatados con la clase real:
\begin {itemize}
\item La inferencia sobre \texttt{Horse} está casi empatada con \texttt{Horse}(35.43\%) y
\texttt{Truck} (37.13\%).
\item La inferencia sobre \texttt{Dog} se confunde con \texttt{Frog}(32.13\%), pero reparte
bastante confianza entre \texttt{Deer}(17.14\%), \texttt{Dog}(18.70\%) y
\texttt{Horse}(14.99\%).
\item La inferencia sobre \texttt{Deer} acierta aunque con un margen muy pequeño frente a
\texttt{Plane}(32.07\%).
\end{itemize}

Estos repartos de predicción apuntan a un modelo pequeño, cuantizado y limitado, no a una salida
rota del acelerador. Esta interpretación encaja con el efecto habitual de la cuantización: reducir
precisión y memoria hace viable la inferencia en hardware pequeño, pero puede estrechar el margen
entre clases cuando el modelo no tiene demasiada capacidad \citep{Jacob_2018_CVPR}.

Hay clases que el modelo reconoce con mucha seguridad, esto nos indica que hay patrones aprendidos
que se conservan correctamente en la \ac{FPGA}.

Aun así, esta prueba no es decisiva para validar la precisión global del modelo porque son solo
10 imágenes. \ac{CIFAR}-10 incluye 60.000 imágenes de 32x32 píxeles repartidas en 10
clases, con 50.000
imágenes de entrenamiento y 10.000 imágenes de test \citep{cifar10Dataset}. Por tanto, una
validación completa debería ejecutar el conjunto de test, o al menos una muestra mucho más amplia y
balanceada por clase. Esa evaluación queda fuera del alcance de este trabajo, donde se ha priorizado
cerrar el flujo completo de inferencia en placa.

Además de la prueba con imágenes sacadas del \textit{dataset}, se ha podido probar la
inferencia sobre
diferentes imágenes descargadas de internet transformándolas mediante la herramienta
\texttt{image\_file\_to\_h.py}.
Teniendo en cuenta que la imagen se debe redimensionar a 32x32 píxeles, si queremos obtener una
predicción correcta con una confianza alta, es importante revisar la imagen reducida y comprobar
que el objeto no se difumina con el entorno, es suficientemente grande en la imagen o no hay
ningún otro objeto en la imagen.

\section{Benchmark de inferencia}

El objetivo principal de rendimiento era comprobar si la \ac{FPGA} reducía la latencia de
la inferencia
frente a una ejecución software en el \ac{ARM} Cortex-A9. Para ello se midieron varias
ejecuciones sobre
la misma imagen y con el mismo temporizador hardware de la plataforma. En esta comparación se usó
el mismo modelo generado por hls4ml en dos caminos distintos: por un lado, el \ac{IP} sintetizado e
integrado en la \ac{FPGA}. Por el otro, el \textit{core} \texttt{C++} generado por hls4ml
ejecutado directamente sobre el
\ac{ARM} \citep{hls4mlDocs}.

La tabla~\ref{tab:benchmarkCifar10FpgaArm} muestra los tiempos obtenidos.

\begin{table}[h!]
\centering
\begin{tabular}{|l|r|r|}
\hline
\textbf{Ruta de inferencia} & \textbf{Tiempo medio} & \textbf{Mejor tiempo} \\ \hline
\ac{FPGA} & 2012 $\mu$s & 2008 $\mu$s \\ \hline
\ac{CPU} (\ac{ARM}) & 171215 $\mu$s & 171161 $\mu$s \\ \hline
\end{tabular}
\caption{Benchmark de inferencia CIFAR-10 entre acelerador FPGA y \textit{core}
\texttt{C++} de hls4ml en ARM}
\label{tab:benchmarkCifar10FpgaArm}
\end{table}

Con estos datos, la aceleración media observada es:

\begin{equation}
\mathrm{speedup} = \frac{171215}{2012} \approx 85.09
\end{equation}

Es decir, la inferencia acelerada en \ac{FPGA} fue aproximadamente 85 veces más rápida
que la ejecución
del \textit{core} \texttt{C++} de hls4ml en el \ac{ARM}. Este resultado es coherente con
la naturaleza de cada ruta. En la
\ac{FPGA}, el modelo se convierte en hardware dedicado y puede explotar paralelismo a nivel de
operaciones y de flujo de datos. En cambio, en el \ac{ARM} se ejecuta una descripción
\texttt{C++} pensada para
ser sintetizada por \ac{HLS}, no una librería de inferencia optimizada para \ac{CPU}
\citep{sze2017efficientprocessingdeepneural}.

Es importante destacar que el tiempo de la ruta \ac{FPGA} incluye el uso real del
sistema: preparación
del \textit{buffer} de salida, coherencia de caché, transferencia \ac{DMA}, ejecución del
\ac{IP} \ac{HLS} y lectura de los
resultados. Por tanto, no se está midiendo únicamente el núcleo matemático del acelerador, sino la
inferencia completa desde el punto de vista de la aplicación \textit{bare-metal}. Aun
incluyendo este coste
de comunicación, el tiempo se mantiene alrededor de los 2 ms.

La ruta \ac{ARM}, por su parte, queda penalizada por varios motivos. El \textit{core}
generado por hls4ml utiliza
tipos de dato de precisión fija, estructuras tipo \texttt{\ac{HLS}::stream} y plantillas
\texttt{C++} orientadas a
describir hardware. Todas estas construcciones tienen sentido cuando el código se va a transformar
en un circuito, pero no son la forma más eficiente de ejecutar una red neuronal en un Cortex-A9.
Por esta razón, el resultado no debe interpretarse como una comparación contra la mejor
implementación posible en \ac{CPU}, sino como una comparación entre el acelerador hardware y la
ejecución software directa del mismo flujo hls4ml.

Si la comparación se realizase con una implementación \ac{SW} optimizada, el tiempo de
ejecución \ac{CPU}
se vería reducido de forma apreciable.
Podemos intentar deducir la mejora de rendimiento que tendría la ejecución \ac{SW} tomando como
referencia \citep{lai2018cmsisnnefficientneuralnetwork}, donde muestran una mejora de rendimiento
de hasta 4.6 veces utilizando un kernel CMSIS-NN \citep{cmsisnnRepo} en una \ac{CNN} entrenada
en el \textit{dataset} \ac{CIFAR}-10, de la misma manera que se hace en este trabajo.
Estos resultados no pueden trasladarse de manera directa a este proyecto, ya que el procesador y
la arquitectura del modelo son distintos.
Sin embargo, sí permiten estimar hasta qué punto es la mejora de una implementación
\ac{SW} optimizada
para un caso de ejecución similar.
Extrapolando esos datos, el tiempo de ejecución de nuestra inferencia en \ac{SW} pasaría a
ser de aproximadamente 37000$\mu$s, todavía por encima de los aproximadamente 2000$\mu$s medidos
en la \ac{FPGA}, que resultaría en una mejora de rendimiento de aproximadamente 18.5 veces.
Por lo tanto, la implementación hardware seguiría manteniendo una ventaja clara respecto a una
ejecución en \ac{CPU}.

Además del tiempo de ejecución, también se observaron diferencias en los valores de salida.

\begin{table}[h!]
\centering
\begin{tabular}{|l|r|r|r|r|}
\hline
\textbf{Imagen} & \textbf{Pred. FPGA} & \textbf{Conf. FPGA} & \textbf{Pred. ARM} &
\textbf{Conf. ARM} \\ \hline
\textbf{Plane} & \textcolor{green!60!black}{Plane} & 71.17\% &
\textcolor{green!60!black}{Plane} & 51.63\% \\ \hline
\textbf{Auto}  & \textcolor{green!60!black}{Auto}  & 60.70\% &
\textcolor{green!60!black}{Auto}  & 46.25\% \\ \hline
\textbf{Bird}  & \textcolor{red}{Truck} & 30.24\% & \textcolor{green!60!black}{Bird}  &
47.21\% \\ \hline
\textbf{Cat}   & \textcolor{green!60!black}{Cat}   & 53.14\% &
\textcolor{green!60!black}{Cat}   & 42.14\% \\ \hline
\textbf{Deer}  & \textcolor{green!60!black}{Deer}  & 34.68\% & \textcolor{red}{Plane} &
30.64\% \\ \hline
\textbf{Dog}   & \textcolor{red}{Frog}  & 32.13\% & \textcolor{red}{Frog}  & 29.04\% \\ \hline
\textbf{Frog}  & \textcolor{green!60!black}{Frog}  & 83.11\% &
\textcolor{green!60!black}{Frog}  & 91.44\% \\ \hline
\textbf{Horse} & \textcolor{red}{Truck} & 37.13\% & \textcolor{green!60!black}{Horse} &
65.24\% \\ \hline
\textbf{Ship}  & \textcolor{red}{Auto}  & 71.26\% & \textcolor{red}{Auto}  & 89.17\% \\ \hline
\textbf{Truck} & \textcolor{green!60!black}{Truck} & 89.08\% &
\textcolor{green!60!black}{Truck} & 62.78\% \\ \hline
\end{tabular}
\caption{Comparación de salida entre la inferencia en FPGA y la ejecución hls4ml en ARM}
\label{tab:comparacionFpgaArmCifar10}
\end{table}

La tabla~\ref{tab:comparacionFpgaArmCifar10} muestra que las dos rutas no producen exactamente
el mismo comportamiento, aunque parten del mismo modelo entrenado.
En las 10 imágenes evaluadas, la inferencia en \ac{FPGA} acierta seis casos mientras que la
ejecución en \ac{ARM} acierta siete.
Estas diferencias deben interpretarse teniendo en cuenta que el modelo final está muy condicionado
por los recursos disponibles en la Zynq Z-7010.

Hay que recordar que la confianza mostrada por terminal no es una salida directa del modelo
hardware. El \ac{IP} devuelve \textit{logits} en Q16.16, y la aplicación calcula después
una \textit{softmax} aproximada
para facilitar la lectura. Por tanto, una diferencia moderada entre \textit{logits} puede
transformarse en
una diferencia grande en el porcentaje final. Este efecto es habitual en clasificadores neuronales:
la \textit{softmax} amplifica la distancia relativa entre las puntuaciones \citep{pmlr-v70-guo17a}.

Las diferencias entre \ac{FPGA} y \ac{ARM} también son razonables desde el punto de vista
numérico. Aunque
ambas rutas parten del mismo modelo entrenado, no ejecutan exactamente la misma implementación
aritmética. La \ac{FPGA} trabaja con el hardware generado por Vitis \ac{HLS}, con los
anchos de palabra,
redondeos y comportamientos de desbordamiento definidos durante la síntesis. El \ac{ARM}, en cambio,
ejecuta una versión \texttt{C++} del modelo mediante tipos software que emulan precisión
fija. Pequeñas
diferencias de redondeo, truncado, orden de acumulación o conversión entre tipos pueden propagarse
a través de las capas convolucionales y modificar la distancia entre los \textit{logits} finales.

Por tanto, la conclusión principal del \textit{benchmark} es doble. En rendimiento, la
\ac{FPGA} ofrece una
mejora muy clara frente a ejecutar el \textit{core} hls4ml en el \ac{ARM}. En precisión
funcional, ambas rutas
mantienen predicciones similares para las imagenes evaluadas, aunque no producen exactamente los
mismos \textit{logits} ni la misma confianza.

\section{Valoración global}

Los resultados muestran que el sistema heterogéneo funciona de extremo a extremo: el
\ac{ARM} prepara
la imagen, la envía a la \ac{FPGA} mediante \ac{AXI} \ac{DMA}, el \ac{IP} \ac{HLS}
ejecuta la inferencia y la aplicación
recupera los \textit{logits} para obtener la clase final. El diseño no se queda en una
simulación aislada,
sino que se ejecuta sobre la placa real.

La mejora temporal frente a ejecutar el \textit{core} hls4ml directamente en el \ac{ARM}
es muy clara, aunque no
debe separarse del contexto de la placa utilizada. La placa Zybo tiene recursos muy ajustados para
una \ac{CNN}, por lo que el modelo final tuvo que ser pequeño y cuidadosamente
cuantizado. A pesar de
ello, conseguir una inferencia en torno a 2 ms en la \ac{FPGA} confirma que el uso de
hardware dedicado
es una vía efectiva para reducir latencia en sistemas empotrados.

La evaluación también deja una enseñanza importante: en proyectos \ac{HLS} no basta con
que el modelo
sea correcto a nivel de código \texttt{C++}. La versión realmente desplegada es el
hardware generado, y por
eso hay que analizar recursos, temporización, comunicación con memoria y comportamiento numérico
una vez integrado en la placa. Esta parte es una pieza clave para poder confiar en los resultados
del sistema completo.
